\documentclass[11pt]{article}
\usepackage[margin=1in]{geometry}
\usepackage{xcolor}
\usepackage[breakable,listings]{tcolorbox}
\tcbuselibrary{listings}
\usepackage{hyperref}
\usepackage{enumitem}
\usepackage{listings}
\usepackage{verbatim}
\usepackage{amsmath}
\usepackage{graphicx}
\usepackage{booktabs}
\usepackage{float}

% Define colors
\definecolor{osaorange}{RGB}{220,115,38}
\definecolor{codebackground}{RGB}{248,248,248}

% Configure hyperref
\hypersetup{
    colorlinks=true,
    linkcolor=blue,
    urlcolor=blue,
    citecolor=blue
}

% Configure listings for code display
\lstset{
  basicstyle=\small\ttfamily,
  breaklines=true,
  breakatwhitespace=true,
  postbreak=\mbox{\textcolor{gray}{$\hookrightarrow$}\space},
  columns=flexible,
  keepspaces=true,
  showstringspaces=false
}

% Code environment using tcolorbox with listings
\newtcblisting{rcode}{
  colback=white,
  colframe=gray!50,
  boxrule=0.5pt,
  arc=0pt,
  left=3pt,
  right=3pt,
  top=3pt,
  bottom=3pt,
  width=\textwidth,
  breakable,
  listing only,
  listing options={
    basicstyle=\small\ttfamily,
    breaklines=true,
    breakatwhitespace=true,
    postbreak=\mbox{\textcolor{gray}{$\hookrightarrow$}\space},
    columns=flexible,
    keepspaces=true,
    showstringspaces=false
  }
}

\title{}
\author{}
\date{}

\begin{document}

% Header box
\begin{tcolorbox}[colback=osaorange,colframe=black,boxrule=2pt,arc=0pt,left=3pt,right=3pt,top=6pt,bottom=6pt,boxsep=2pt]
\centering
\textcolor{white}{\textbf{\Large H524 Introduction to Biostatistics}}\\
\textcolor{white}{\textbf{\Large Week 5 Lab: Hypothesis Testing}}\\
\textcolor{white}{\textbf{\Large Fall 2025}}
\end{tcolorbox}

\vspace{8pt}

\noindent\textbf{Format:} Online\\
\textbf{Tools Required:} R, RStudio\\
\textbf{Estimated Time:} 90 minutes

\section*{Problem 1: Two-Sided t-Test (Body Temperature)}

\textbf{Research Question:} Is the true mean body temperature different from the commonly cited value of 98.6°F?

\textbf{Data:} Body temperatures from $n = 15$ healthy adults:
\begin{rcode}
temps <- c(98.4, 98.0, 98.2, 98.6, 98.5, 98.3, 98.1, 98.7,
           98.2, 98.4, 98.3, 98.0, 98.5, 98.4, 98.2)
\end{rcode}

\subsection*{Part (a): State the hypotheses}

\textbf{Null hypothesis:}

\vspace{1cm}

\textbf{Alternative hypothesis:}

\vspace{1cm}

\textbf{Note:} Is this a one-sided or two-sided test? Why?

\vspace{1cm}

\subsection*{Part (b): Calculate descriptive statistics}

\textbf{R code:}
\begin{rcode}
n <- length(temps)
xbar <- mean(temps)
s <- sd(temps)

n
xbar
s
\end{rcode}

\textbf{Results:}
\begin{itemize}
\item Sample size: $n = $
\item Sample mean: $\bar{x} = $
\item Sample standard deviation: $s = $
\end{itemize}

\vspace{1cm}

\subsection*{Part (c): Conduct the hypothesis test}

\textbf{Step 1: Calculate the standard error}

\vspace{2cm}

\textbf{Step 2: Calculate the t-statistic}

\vspace{2cm}

\textbf{Step 3: Determine degrees of freedom}

\vspace{1cm}

\textbf{Step 4: Calculate the p-value}

For a two-sided test, we calculate the probability in both tails:

\vspace{1cm}

\textbf{R verification:}
\begin{rcode}
# Manual calculation
SE <- s / sqrt(n)
t_stat <- (xbar - 98.6) / SE
p_value <- 2 * pt(abs(t_stat), df = n - 1, lower.tail = FALSE)

SE
t_stat
p_value

# Using t.test function
result <- t.test(temps, mu = 98.6, alternative = "two.sided")
result
\end{rcode}

\vspace{1cm}

\subsection*{Part (d): Make a decision and interpret}

\textbf{Decision at $\alpha = 0.05$:}

\vspace{1cm}

\textbf{Interpretation:}

\vspace{3cm}

\textbf{95\% Confidence Interval:}

\vspace{1cm}

\pagebreak

\section*{Problem 2: One-Sided t-Test (Blood Pressure)}

\textbf{Research Question:} Does a new blood pressure medication lower mean systolic blood pressure below 140 mmHg?

\textbf{Data:} Systolic blood pressure from $n = 20$ patients after treatment:
\begin{rcode}
bp <- c(138, 135, 142, 128, 145, 137, 132, 140, 136, 134,
        139, 141, 133, 130, 143, 135, 138, 142, 136, 139)
\end{rcode}

\subsection*{Part (a): State the hypotheses}

\textbf{Null hypothesis:}

\vspace{1cm}

\textbf{Alternative hypothesis:}

\vspace{1cm}

\textbf{Note:} Is this a one-sided or two-sided test? Why?

\vspace{1cm}

\subsection*{Part (b): Calculate descriptive statistics}

\textbf{R code:}
\begin{rcode}
n <- length(bp)
xbar <- mean(bp)
s <- sd(bp)

n
xbar
s
\end{rcode}

\textbf{Results:}
\begin{itemize}
\item Sample size: $n = $
\item Sample mean: $\bar{x} = $
\item Sample standard deviation: $s = $
\end{itemize}

\vspace{1cm}

\subsection*{Part (c): Conduct the hypothesis test}

\textbf{Step 1: Calculate the standard error}

\vspace{2cm}

\textbf{Step 2: Calculate the t-statistic}

\vspace{2cm}

\textbf{Step 3: Determine degrees of freedom}

\vspace{1cm}

\textbf{Step 4: Calculate the p-value (one-sided, lower tail)}

\vspace{1cm}

\textbf{R verification:}
\begin{rcode}
# Manual calculation
SE <- s / sqrt(n)
t_stat <- (xbar - 140) / SE
p_value <- pt(t_stat, df = n - 1)

SE
t_stat
p_value

# Using t.test function (one-sided)
result <- t.test(bp, mu = 140, alternative = "less")
result
\end{rcode}

\vspace{1cm}

\subsection*{Part (d): Make a decision and interpret}

\textbf{Decision at $\alpha = 0.05$:}

\vspace{1cm}

\textbf{Interpretation:}

\vspace{3cm}

\textbf{Clinical significance:}

\vspace{2cm}

\pagebreak

\section*{Problem 3: Type I and Type II Errors}

Using the blood pressure example from Problem 2, explain the concepts of Type I and Type II errors.

\subsection*{Type I Error ($\alpha$)}

\textbf{Definition:}

\vspace{1cm}

\textbf{In this context:}

\vspace{1cm}

\textbf{Probability:}

\vspace{1cm}

\textbf{Consequence:}

\vspace{3cm}

\textbf{Control:}

\vspace{1cm}

\subsection*{Type II Error ($\beta$)}

\textbf{Definition:}

\vspace{1cm}

\textbf{In this context:}

\vspace{1cm}

\textbf{Probability:}

\vspace{1cm}

\textbf{Consequence:}

\vspace{3cm}

\textbf{Control:}

\vspace{2cm}

\subsection*{Statistical Power}

\textbf{Definition:}

\vspace{1cm}

\textbf{Interpretation:}

\vspace{1cm}

\textbf{Typical goal:}

\vspace{1cm}

\subsection*{Trade-off}

There is an inverse relationship between Type I and Type II errors. Explain:

\vspace{3cm}

\pagebreak

\section*{Problem 4: Power Analysis}

Understanding statistical power and sample size determination.

\subsection*{Part (a): Calculate power for a given sample size}

\textbf{Scenario:} We want to detect a medium effect (Cohen's $d = 0.5$) with $n = 30$ subjects at $\alpha = 0.05$ (two-sided test).

\textbf{R code:}
\begin{rcode}
library(pwr)

power_result <- pwr.t.test(n = 30,
                            d = 0.5,
                            sig.level = 0.05,
                            type = "one.sample",
                            alternative = "two.sided")
power_result
\end{rcode}

\textbf{Result:} Power $\approx$

\vspace{1cm}

\textbf{Interpretation:}

\vspace{2cm}

\textbf{Assessment:}

\vspace{2cm}

\subsection*{Part (b): Calculate sample size for desired power}

\textbf{Question:} How many subjects do we need to achieve 80\% power for detecting a medium effect ($d = 0.5$)?

\textbf{R code:}
\begin{rcode}
n_result <- pwr.t.test(d = 0.5,
                        power = 0.80,
                        sig.level = 0.05,
                        type = "one.sample",
                        alternative = "two.sided")
n_result
\end{rcode}

\textbf{Result:} Required $n = $

\vspace{1cm}

\textbf{Interpretation:}

\vspace{2cm}

\subsection*{Part (c): Factors affecting power}

Power increases when:

\vspace{4cm}

\textbf{Key insight:}

\vspace{2cm}

\pagebreak

\section*{Problem 5: Relationship Between Confidence Intervals and Hypothesis Tests}

There is a fundamental equivalence between confidence intervals and hypothesis tests.

\subsection*{Equivalence Principle}

For a two-sided test at significance level $\alpha$:

\textbf{If the $(1-\alpha) \times 100\%$ confidence interval contains the null value $\mu_0$:}

\vspace{2cm}

\textbf{If the $(1-\alpha) \times 100\%$ confidence interval does NOT contain $\mu_0$:}

\vspace{2cm}

\subsection*{Demonstration with Body Temperature Data}

Using the body temperature data from Problem 1:

\textbf{95\% Confidence Interval:}

\vspace{1cm}

\textbf{Null value:} $\mu_0 = 98.6$ °F

\textbf{Question:} Is 98.6 inside the interval?

\vspace{1cm}

\textbf{Answer:}

\vspace{1cm}

\textbf{Conclusion from CI:}

\vspace{2cm}

\textbf{Conclusion from hypothesis test:}

\vspace{2cm}

\textbf{Result:} Do both methods agree?

\vspace{1cm}

\subsection*{Advantages of Confidence Intervals}

Confidence intervals provide MORE information than hypothesis tests:

\vspace{5cm}

\textbf{Recommendation:}

\vspace{1cm}

\pagebreak

\section*{Problem 6: Checking the Normality Assumption}

The t-test assumes that the data come from a normally distributed population. We should check this assumption.

\subsection*{Methods for Checking Normality}

\textbf{1. Visual Methods:}

\vspace{3cm}

\textbf{2. Formal Test:}

\vspace{1cm}

\subsection*{Applying to Body Temperature Data}

\textbf{R code for diagnostic plots:}
\begin{rcode}
# Create diagnostic plots
par(mfrow = c(1, 3))

# Histogram
hist(temps, breaks = 5, main = "Histogram of Body Temps",
     xlab = "Temperature (F)", col = "lightblue", border = "white")

# Boxplot
boxplot(temps, main = "Boxplot", ylab = "Temperature (F)")

# Q-Q plot (most important for normality)
qqnorm(temps, main = "Q-Q Plot", pch = 19, col = "blue")
qqline(temps, col = "red", lwd = 2)

par(mfrow = c(1, 1))
\end{rcode}

\vspace{1cm}

\subsection*{Shapiro-Wilk Test}

\textbf{R code:}
\begin{rcode}
shapiro.test(temps)
\end{rcode}

\textbf{Hypotheses:}
\begin{itemize}
\item $H_0$:
\item $H_A$:
\end{itemize}

\vspace{1cm}

\textbf{Results:}
\begin{itemize}
\item Test statistic: $W = $
\item P-value: $p = $
\end{itemize}

\vspace{1cm}

\textbf{Decision:}

\vspace{1cm}

\textbf{Interpretation:}

\vspace{2cm}

\subsection*{Assessment}

\textbf{Visual assessment:}

\vspace{3cm}

\textbf{Formal test:}

\vspace{1cm}

\textbf{Conclusion:}

\vspace{2cm}

\subsection*{What if normality is violated?}

If normality is violated:

\vspace{5cm}

\pagebreak

\section*{Summary of Key Concepts}

\subsection*{Steps for Hypothesis Testing}

\begin{enumerate}
\item \textbf{State hypotheses} ---
\item \textbf{Choose significance level} ---
\item \textbf{Check assumptions} ---
\item \textbf{Calculate test statistic} ---
\item \textbf{Find p-value} ---
\item \textbf{Make decision} ---
\item \textbf{Interpret in context} ---
\end{enumerate}

\subsection*{One-Sample t-Test Formula}

\textbf{Test statistic:}
$$t = \frac{\bar{x} - \mu_0}{s / \sqrt{n}} \quad \text{with } df = n - 1$$

\textbf{Two-sided p-value:}
$$p = 2 \times P(T \geq |t|)$$

\textbf{One-sided p-value (lower tail):}
$$p = P(T \leq t)$$

\textbf{One-sided p-value (upper tail):}
$$p = P(T \geq t)$$

\subsection*{Error Types}

\begin{center}
\begin{tabular}{lcc}
\toprule
& \textbf{$H_0$ True} & \textbf{$H_0$ False} \\
\midrule
\textbf{Reject $H_0$} & & \\
\textbf{Fail to Reject $H_0$} & & \\
\bottomrule
\end{tabular}
\end{center}

\subsection*{Important Relationships}

\begin{itemize}
\item \textbf{CI and hypothesis test:}
\item \textbf{Sample size and power:}
\item \textbf{Effect size and power:}
\item \textbf{Alpha and power:}
\end{itemize}

\vspace{2cm}

\noindent\textit{Week 5 Lab: Hypothesis Testing}\\
\noindent\textit{Fall 2025}

\end{document}
